\chapter{Random Walks on Networks and Groups}

\begin{center}\itshape
A random walker on a graph and an electric current in a network of resistors solve the same equations.\\
So a question of probability, does the walker come back?, becomes a question of physics:\\
is the resistance to infinity finite? And physics has excellent tools.
\end{center}

In 1847 Gustav Kirchhoff, then a twenty-three year old \emph{Privatdozent} in Berlin, published the solution of a problem he had posed as a student: given a network of wires and a battery, find the
current in each wire. His two laws, that the currents entering a node sum to zero and that the voltage drops around a loop sum to zero, determine everything, and the voltages solve a discrete Laplace
equation. A century later probabilists noticed that the same equation governs a random walk on the network: the probability that the walker, started at $x$, reaches $a$ before $b$ is the voltage at $x$ when
$a$ is held at one volt and $b$ is earthed (Volume~3, Exercise~7.8).

The dictionary was made systematic by Crispin Nash-Williams, in a 1959 paper in the \emph{Proceedings of the Cambridge Philosophical Society}, and popularised by Peter Doyle and Laurie Snell in their little
book \emph{Random Walks and Electric Networks} (1984), which proves P\'olya's theorem again by cutting and shorting wires. Its key physical principle is
older: Lord Rayleigh knew that increasing the resistance of a single wire cannot decrease the resistance of the network, and William Thomson (Lord Kelvin) that the current chooses, among all the flows with
the same input, the one that dissipates the least energy. Terry Lyons turned Thomson's principle into a criterion for transience in 1983: a walk escapes to infinity exactly when one can send a unit of
current to infinity with finite energy.

Groups enter through their Cayley graphs. In his 1958 Cornell thesis, written under Mark Kac, Harry Kesten asked how fast the return probability of a random walk on a group decays, and found that the
answer, a number between $0$ and $1$, detects a deep property of the group, its \emph{amenability}. For the free group he computed it exactly: $\frac{\sqrt{2k-1}}k$, the number behind the Ramanujan graphs of the
next chapter.

\begin{reminder}[from Volume 3 and Appendix A]
A walk returns to its starting point infinitely often (\emph{recurrence}) if and only if the expected number of returns $\sum_nu_n$ is infinite (Volume~3, Theorem~7.12), and simple random walk on $\Z^d$ is recurrent for
$d\le2$ and transient for $d\ge3$ (P\'olya, Volume~3, Theorem~7.13). The discrete Dirichlet problem (Volume~3, Theorem~7.14). Rayleigh's quotient and the Laplacian matrix of a graph (Volume~11, Appendix~A). The Catalan numbers
$C_m=\frac1{m+1}\binom{2m}m$ count the paths of $2m$ steps $\pm1$ from $0$ to $0$ that never go below $0$.
\end{reminder}

\section{Networks and reversible walks}

A \emph{network} is a connected graph $G=(V,E)$, finite or countable and locally finite, with a \emph{conductance} $c(e)>0$ on each edge; $r(e)=1/c(e)$ is its \emph{resistance}. We write $c(xy)$ for the conductance of the edge
$\{x,y\}$ ($0$ if there is none), $\pi(x)=\sum_yc(xy)$, and define the \emph{network walk} by
\[
P(x,y)=\frac{c(xy)}{\pi(x)} .
\]
With all conductances $1$ it is simple random walk. The walk is \emph{reversible}: $\pi(x)P(x,y)=c(xy)=\pi(y)P(y,x)$. Conversely every irreducible Markov chain with a positive measure $\pi$ satisfying this
\emph{detailed balance} is a network walk, with $c(xy)=\pi(x)P(x,y)$. On a finite network, $\pi/\pi(V)$ is a stationary law.

A function $f$ on $V$ is \emph{harmonic} at $x$ if $f(x)=\sum_yP(x,y)f(y)$, that is if $\Delta f(x)=0$, where
\[
\Delta f(x)=\sum_yc(xy)\bigl(f(x)-f(y)\bigr)=\pi(x)\Bigl(f(x)-\sum_yP(x,y)f(y)\Bigr).
\]
(With this sign, $\Delta$ is the \emph{positive} Laplacian of Appendix~A.) The \emph{energy} of $f$ is $\mathcal E(f)=\frac12\sum_{x,y}c(xy)(f(x)-f(y))^2$, and for finitely supported $f$ or finite $V$ summation by parts gives
\begin{equation}\label{eq:green}
\begin{aligned}
&\sum_xg(x)\Delta f(x)\\
&\qquad=\frac12\sum_{x,y}c(xy)\bigl(f(x)-f(y)\bigr)\bigl(g(x)-g(y)\bigr).
\end{aligned}
\end{equation}

\begin{lemma}[maximum principle]\label{lem:maxprinciple}
Let $G$ be finite, $Z\subseteq V$ non-empty, and $f$ harmonic on $V\setminus Z$. Then $\max_Vf=\max_Zf$. Consequently, for every $g$ on $Z$ there is exactly one $f$ harmonic on $V\setminus Z$ with $f=g$ on $Z$, namely
$f(x)=\mathbb E_x[g(X_{\tau_Z})]$, where $\tau_Z$ is the hitting time of $Z$.
\end{lemma}

\begin{proof}
If the maximum $M$ is attained at $x\notin Z$, it is attained at every neighbour of $x$ (an average equal to the maximum), hence along a path from $x$ to $Z$, by connectedness. Uniqueness follows by applying this to
$\pm(f_1-f_2)$. The function $\mathbb E_x[g(X_{\tau_Z})]$ equals $g$ on $Z$ and is harmonic off $Z$ by conditioning on the first step; $\tau_Z<\infty$ almost surely since from every vertex there is a path to $Z$ of length at most
$|V|$, followed with probability at least some $\delta>0$, so $P_x(\tau_Z>m|V|)\le(1-\delta)^m$.
\end{proof}

\section{Voltage, current and effective resistance}

Fix a finite network, a vertex $a$ and a set $Z\not\ni a$. The \emph{voltage} is the harmonic function on $V\setminus(\{a\}\cup Z)$ with $v(a)=1$ and $v=0$ on $Z$: by Lemma~\ref{lem:maxprinciple},
\[
v(x)=P_x(\tau_a<\tau_Z) .
\]
The \emph{current} is $i(xy)=c(xy)(v(x)-v(y))$ (Ohm's law), antisymmetric, with $\sum_yi(xy)=\Delta v(x)=0$ at every $x\notin\{a\}\cup Z$ (Kirchhoff's node law). The total current leaving $a$ is
$I=\Delta v(a)=\sum_yi(ay)$, and the \emph{effective resistance} between $a$ and $Z$ is $R(a\leftrightarrow Z)=1/I$: the voltage needed to push a unit current.

\begin{proposition}[escape probability]\label{prop:escape}
$P_a(\tau_Z<\tau_a^+)=\dfrac1{\pi(a)R(a\leftrightarrow Z)}$, where $\tau_a^+=\min\{n\ge1:X_n=a\}$.
\end{proposition}

\begin{proof}
Conditioning on the first step, $P_a(\tau_a^+<\tau_Z)=\sum_yP(a,y)v(y)$, so $P_a(\tau_Z<\tau_a^+)=\sum_yP(a,y)(v(a)-v(y))=\frac{\Delta v(a)}{\pi(a)}=\frac I{\pi(a)}$.
\end{proof}

A \emph{flow} from $a$ to $Z$ is an antisymmetric function $\theta$ on oriented edges with \emph{divergence} $\operatorname{div}\theta(x)=\sum_y\theta(xy)$ equal to $0$ for $x\notin\{a\}\cup Z$; it is a \emph{unit} flow if $\operatorname{div}\theta(a)=1$. Its energy is
$\mathcal E(\theta)=\frac12\sum_{x,y}r(xy)\theta(xy)^2$. The \emph{unit current} is $i/I$.

\begin{theorem}[Thomson's principle]\label{thm:thomson}
We have
\[
R(a\leftrightarrow Z)=\min\{\mathcal E(\theta):\theta\text{ a unit flow from }a\text{ to }Z\},
\]
and the minimum is attained exactly at the unit current.
\end{theorem}

\begin{proof}
Let $j=i/I$, with potential $u=v/I$, so that $r(xy)j(xy)=u(x)-u(y)$. For any antisymmetric $\eta$, summation by parts gives
\begin{equation}\label{eq:sumparts}
\begin{aligned}
\frac12\sum_{x,y}r(xy)j(xy)\eta(xy)&=\frac12\sum_{x,y}(u(x)-u(y))\eta(xy)\\
&=\sum_xu(x)\operatorname{div}\eta(x).
\end{aligned}
\end{equation}
With $\eta=j$: $\mathcal E(j)=\sum_xu(x)\operatorname{div}j(x)=u(a)\cdot1=\frac1I$, since $u=0$ on $Z$ and $\operatorname{div}j=0$ elsewhere. If $\theta$ is a unit flow, $\eta=\theta-j$ has divergence $0$ off $Z$, so \eqref{eq:sumparts} vanishes and
$\mathcal E(\theta)=\mathcal E(j)+\mathcal E(\theta-j)\ge\mathcal E(j)$, with equality iff $\theta=j$.
\end{proof}

\begin{corollary}[Rayleigh's monotonicity law]\label{cor:rayleigh}
If the conductances are increased (in particular if edges are added, or vertices are \emph{shorted}, that is glued together), the effective resistance $R(a\leftrightarrow Z)$ does not increase.
\end{corollary}

\begin{proof}
Each energy $\mathcal E(\theta)=\frac12\sum r\theta^2$ decreases when the resistances decrease, and so does their minimum. Adding an edge is increasing its conductance from $0$ (a unit flow of the old network is one of the new
network with $\theta=0$ on the new edge). Shorting $x$ and $y$ is the limit of adding an edge of conductance $n\to\infty$ between them; effective resistances of finite networks are continuous in the conductances (the
voltage solves a linear system with an invertible matrix).
\end{proof}

The \emph{Dirichlet principle} is the dual statement, with functions instead of flows (Exercise~\ref{ex:dirichletprinciple}): $\frac1{R(a\leftrightarrow Z)}=\min\{\mathcal E(f):f(a)=1,\ f=0\text{ on }Z\}$.

\begin{theorem}[commute time identity; Chandra, Raghavan, Ruzzo, Smolensky and Tiwari, 1989]\label{thm:commute}
On a finite network, with $c(G)=\sum_x\pi(x)=2\sum_ec(e)$, for all vertices $a\ne b$,
\[
\mathbb E_a\tau_b+\mathbb E_b\tau_a=c(G)\,R(a\leftrightarrow b) .
\]
\end{theorem}

\begin{proof}
Let $\varphi_b(x)=\mathbb E_x\tau_b$, finite as in Lemma~\ref{lem:maxprinciple}. Conditioning on the first step, $\varphi_b(x)=1+\sum_yP(x,y)\varphi_b(y)$ for $x\ne b$, that is $\Delta\varphi_b(x)=\pi(x)$. By \eqref{eq:green} with
$g=1$, $\sum_x\Delta\varphi_b(x)=0$, so $\Delta\varphi_b(b)=\pi(b)-c(G)$. The same holds for $\varphi_a$. Hence $\psi=\varphi_b-\varphi_a$ satisfies $\Delta\psi=0$ off $\{a,b\}$, $\Delta\psi(a)=\pi(a)-(\pi(a)-c(G))=c(G)$ and
$\Delta\psi(b)=-c(G)$.

The potential $u$ of the unit current from $a$ to $b$ satisfies $\Delta u=\delta_a-\delta_b$ and $u(a)-u(b)=R(a\leftrightarrow b)$. The function $\psi/c(G)-u$ is harmonic everywhere, hence constant by the maximum principle (with
$Z$ any single vertex). So $\psi(a)-\psi(b)=c(G)R(a\leftrightarrow b)$, and $\psi(a)-\psi(b)=\varphi_b(a)-0-0+\varphi_a(b)$.
\end{proof}

On the path $0,1,\dots,n$ with unit conductances, $c(G)=2n$ and $R(0\leftrightarrow n)=n$, so $\mathbb E_0\tau_n+\mathbb E_n\tau_0=2n^2$, and by symmetry $\mathbb E_0\tau_n=n^2$: the gambler's ruin time of Volume~3, now from Ohm's law. The proof also
gives Kac's formula $\mathbb E_a\tau_a^+=c(G)/\pi(a)$ (Exercise~\ref{ex:kac}).

\section{Recurrence and transience by electricity}

Let now $G$ be an infinite network, $a$ a vertex, and $G_n$ an increasing sequence of finite connected sets of vertices containing $a$ with union $V$. Let $G_n^*$ be the finite network obtained from $G$ by shorting all the
vertices outside $G_n$ into a single vertex $z_n$ (keeping the edges, now possibly multiple, from $G_n$ to $z_n$). By Rayleigh, $R(a\leftrightarrow z_n)$ in $G_n^*$ is non-decreasing in $n$ (shorting fewer vertices); its limit
$R(a\leftrightarrow\infty)\in]0,\infty]$ is the \emph{resistance to infinity}.

\begin{theorem}\label{thm:recurrentR}
The network walk is recurrent if and only if $R(a\leftrightarrow\infty)=\infty$. More precisely $P_a(\tau_a^+=\infty)=\frac1{\pi(a)R(a\leftrightarrow\infty)}$.
\end{theorem}

\begin{proof}
Until it leaves $G_n$, the walk on $G$ and the walk on $G_n^*$ coincide, so by Proposition~\ref{prop:escape} the probability that the walk on $G$ from $a$ leaves $G_n$ before returning to $a$ is $\frac1{\pi(a)R(a\leftrightarrow z_n)}$. These events
decrease to $\{\tau_a^+=\infty\}$: a walk that never returns to $a$ leaves every finite set, since a walk confined to the finite connected set $G_n$ forever would visit $a$ (Lemma~\ref{lem:maxprinciple}). Recurrence means that
$P_a(\tau_a^+<\infty)=1$ (Volume~3, Theorem~7.12, whose proof applies to any Markov chain).
\end{proof}

A \emph{flow from $a$ to infinity} is an antisymmetric $\theta$ on oriented edges with $\operatorname{div}\theta(x)=0$ for $x\ne a$ and $\operatorname{div}\theta(a)=1$.

\begin{theorem}[T.~Lyons, 1983]\label{thm:lyons}
If there is a flow from $a$ to infinity with finite energy, the network walk is transient.
\end{theorem}

\begin{proof}
Let $\theta$ be such a flow. Its restriction to the edges of $G_n^*$ (the edges leaving $G_n$ now ending at $z_n$) is a unit flow from $a$ to $z_n$, with energy at most $\mathcal E(\theta)$. By Thomson,
$R(a\leftrightarrow z_n)\le\mathcal E(\theta)$ for all $n$, so $R(a\leftrightarrow\infty)<\infty$.
\end{proof}

The converse holds too: when $R(a\leftrightarrow\infty)<\infty$ the unit currents of $G_n^*$ converge to a flow to infinity of energy $R(a\leftrightarrow\infty)$ (Lyons and Peres, Chapter~2). For recurrence one needs lower bounds on
resistance, and the simplest are cuts. A set $\Pi$ of edges is a \emph{cutset} separating $a$ from infinity if every infinite simple path from $a$ uses an edge of $\Pi$.

\begin{theorem}[Nash-Williams, 1959]\label{thm:nashwilliams}
If $\Pi_1,\Pi_2,\dots$ are disjoint finite cutsets separating $a$ from infinity, then
\[
R(a\leftrightarrow\infty)\ge\sum_k\Bigl(\sum_{e\in\Pi_k}c(e)\Bigr)^{-1}.
\]
In particular, if the sum diverges, the walk is recurrent.
\end{theorem}

\begin{proof}
Fix $m$ and $n$ so large that $G_n$ contains both endpoints of every edge of $\Pi_1,\dots,\Pi_m$. A unit flow $\theta$ from $a$ to $z_n$ in $G_n^*$ sends a total of $1$ across each cutset: the set $S$ of vertices reachable from $a$
without using $\Pi_k$ is finite and does not contain $z_n$, so summing the divergence over $S$ gives $1=\sum_{x\in S}\operatorname{div}\theta(x)=\sum\theta(xy)$ over the edges leaving $S$, all in $\Pi_k$. So
$\sum_{e\in\Pi_k}|\theta(e)|\ge1$ and by Cauchy--Schwarz
\[
1\le\Bigl(\sum_{e\in\Pi_k}|\theta(e)|\Bigr)^2\le\sum_{e\in\Pi_k}c(e)\cdot\sum_{e\in\Pi_k}r(e)\theta(e)^2 .
\]
Summing over $k\le m$ and using disjointness, $\mathcal E(\theta)\ge\sum_{k\le m}(\sum_{\Pi_k}c)^{-1}$. Thomson's principle and $n,m\to\infty$ conclude.
\end{proof}

\begin{figure}[ht]
\centering
\fig{ch08-cuts}
\caption{Left: in $\Z^2$ the edges joining the boundary of the box $[-k,k]^2$ to the outside form disjoint cutsets $\Pi_k$, with $8k+4$ edges each; $\sum_k\frac1{8k+4}=\infty$, so the walk is recurrent. Right: on the binary tree, the
flow that splits equally at each vertex sends $2^{-n}$ through each of the $2^n$ edges of generation $n$, with energy $\sum_n2^n4^{-n}<\infty$, so the walk is transient.}
\label{fig:cuts}
\end{figure}

\begin{corollary}\label{cor:z2rec}
Simple random walk on $\Z^2$ is recurrent, and so is simple random walk on every subgraph of $\Z^2$.
\end{corollary}

\begin{proof}
The cutsets $\Pi_k$ of Figure~\ref{fig:cuts} give $R(0\leftrightarrow\infty)\ge\sum_k\frac1{8k+4}=\infty$. Removing edges increases the resistance (Rayleigh).
\end{proof}

\begin{theorem}\label{thm:zdtransient}
Simple random walk on $\Z^d$, $d\ge3$, is transient; so is simple random walk on every graph containing $\Z^d$, $d\ge3$, as a subgraph.
\end{theorem}

\begin{proof}
We build a finite-energy flow from $0$ to infinity out of the electric field of a point charge. Let $\sigma_{d-1}$ be the area of the unit sphere and $F(y)=\frac{y}{\sigma_{d-1}|y|^d}$ on $\R^d\setminus\{0\}$. It is
divergence-free there (a direct computation: $\operatorname{div}(y|y|^{-d})=d|y|^{-d}-d|y|^{-d}=0$), and its flux out of any sphere around $0$ is $1$. For $x\in\Z^d$ let $Q_x=x+[-\frac12,\frac12]^d$, and for neighbours $x,y$ let $\theta(xy)$ be
the flux of $F$ through the common face of $Q_x$ and $Q_y$, oriented from $Q_x$ to $Q_y$. Then $\operatorname{div}\theta(x)$ is the flux of $F$ out of $\partial Q_x$: $0$ for $x\ne0$ by the divergence theorem, and $1$ for $x=0$ (compare with a small
sphere around $0$, by the divergence theorem in the region between them). A face at distance at least $\frac12|x|$ from $0$ (for $|x|\ge2$) carries a flux at most $\sup|F|\le C|x|^{1-d}$, so
\[
\mathcal E(\theta)\le C'+\sum_{x\ne0}2d\,C^2|x|^{2-2d}<\infty ,
\]
since $\sum_{x\in\Z^d\setminus\{0\}}|x|^{-s}<\infty$ for $s>d$, and $2d-2>d$ when $d\ge3$. Theorem~\ref{thm:lyons} applies. For a graph containing $\Z^d$, the same flow (zero on the other edges) works.
\end{proof}

This is P\'olya's theorem again, proved without a single return probability, and much more robust: the triangular lattice, $\Z^2$ with a few extra edges, or any subgraph of $\Z^2$ is recurrent; the cubic lattice with
any edges added, or a slab $\Z^2\times\{0,\dots,k\}$ thickened into $\Z^3$, is transient. (The slab itself is recurrent: Exercise~\ref{ex:slab}.)

\begin{beauty}
The walker never sees the whole lattice, yet whether it comes home depends on a global quantity, the resistance to infinity. In dimension $2$, the cuts at distance $k$ have about $k$ edges and the resistance grows
like $\sum\frac1k$, a logarithm; in dimension $3$ the field of a point charge carries the current away with energy $\sum k^{2}\cdot k^{-4}$, and the walker escapes. The planar Brownian motion of Chapter~5, recurrent
only for neighbourhoods, sits exactly on the logarithmic borderline.
\end{beauty}

\section{Trees and groups}

\begin{proposition}\label{prop:treetransient}
On the $d$-regular tree, $d\ge3$, simple random walk is transient, with $P_a(\tau_a^+=\infty)=\frac{d-2}{d-1}$.
\end{proposition}

\begin{proof}
By symmetry the voltage from $a$ to the vertices at distance $n$ depends only on the distance, so shorting each sphere changes nothing (no current flows within a sphere), and the network becomes a path in which
the $k$-th edge consists of $d(d-1)^{k-1}$ parallel unit resistors. Hence $R(a\leftrightarrow\infty)=\sum_{k\ge1}\frac1{d(d-1)^{k-1}}=\frac{d-1}{d(d-2)}$, and $\pi(a)=d$ in Theorem~\ref{thm:recurrentR}.
\end{proof}

Let $\Gamma$ be a group generated by a finite symmetric set $S$ ($s\in S\Rightarrow s^{-1}\in S$). Its \emph{Cayley graph} has vertex set $\Gamma$ and an edge between $g$ and $gs$ for $s\in S$; simple random walk on it is
$X_n=s_1s_2\cdots s_n$ with $s_i$ independent and uniform on $S$. Let $p_n=P(X_n=e)$, the return probability. Since $p_{2m+2n}\ge p_{2m}p_{2n}$ (return at time $2m$, then again at time $2m+2n$), the limit
\[
\rho=\lim_{n\to\infty}p_{2n}^{1/(2n)}=\sup_np_{2n}^{1/(2n)}
\]
exists by Fekete's lemma (for $a_n=-\log p_{2n}$, which is subadditive). It is the \emph{spectral radius} of the walk. On $\Z^d$, $p_{2n}\asymp n^{-d/2}$ and $\rho=1$.

\begin{lemma}[Schur's test]\label{lem:schur}
Let $P$ be a symmetric non-negative kernel on a countable set and $f>0$ with $\sum_yP(x,y)f(y)\le\lambda f(x)$ for all $x$. Then $|\langle g,Ph\rangle|\le\lambda\|g\|_2\|h\|_2$ for $g,h\in\ell^2$, so $\|P\|_{\ell^2\to\ell^2}\le\lambda$.
\end{lemma}

\begin{proof}
$|g(x)h(y)|\le\frac12\bigl(g(x)^2\frac{f(y)}{f(x)}+h(y)^2\frac{f(x)}{f(y)}\bigr)$. Multiply by $P(x,y)$ and sum: the first term gives at most $\frac12\sum_xg(x)^2\lambda$, and by symmetry the second at most $\frac12\sum_yh(y)^2\lambda$.
\end{proof}

\begin{theorem}[Kesten, 1959]\label{thm:kesten}
For simple random walk on the free group $F_k$, $k\ge1$, with generators $a_1^{\pm1},\dots,a_k^{\pm1}$,
\[
\rho=\frac{\sqrt{2k-1}}k .
\]
\end{theorem}

\begin{proof}
Every element of $F_k$ is a unique reduced word, and the Cayley graph is the $2k$-regular tree. If $|g|$ denotes the length of the reduced word, from $g\ne e$ exactly one of the $2k$ generators shortens the word and the
others lengthen it.

\emph{Upper bound.} Let $f(g)=(2k-1)^{-|g|/2}$. For $g\ne e$, $\sum_sP(g,gs)f(gs)=\frac1{2k}\bigl(\sqrt{2k-1}+\frac{2k-1}{\sqrt{2k-1}}\bigr)f(g)=\frac{\sqrt{2k-1}}kf(g)$, and at $e$ the sum is $(2k-1)^{-1/2}\le\frac{\sqrt{2k-1}}k$. Lemma~\ref{lem:schur}
gives $\|P\|\le\frac{\sqrt{2k-1}}k$, so $p_{2n}=\langle\delta_e,P^{2n}\delta_e\rangle\le\|P\|^{2n}$.

\emph{Lower bound.} The length $|X_n|$ is a Markov chain on $\{0,1,\dots\}$, going from $0$ to $1$, and from $j\ge1$ up with probability $p=\frac{2k-1}{2k}$ and down with probability $q=\frac1{2k}$. A return to $e$ at time $2n$
happens in particular when the length goes up at the first step, then follows one of the $C_{n-1}$ paths from $1$ to $1$ in $2n-2$ steps that stay $\ge1$, then goes down: $p_{2n}\ge C_{n-1}p^{n-1}q^n$. Since
$C_m\ge\frac{4^m}{(2m+1)(m+1)}$ (the central binomial coefficient is the largest of the $2m+1$ coefficients summing to $4^m$), $p_{2n}^{1/(2n)}\to$ at least $\sqrt{4pq}=\frac{\sqrt{2k-1}}k$.
\end{proof}

For $k\ge2$, $\rho<1$: the return probabilities decay exponentially and the walk is transient, more strongly than on any $\Z^d$. For $k=1$, $F_1=\Z$ and $\rho=1$. Kesten proved much more: for every symmetric finitely
supported walk on a group, $\rho=1$ if and only if the group is \emph{amenable} (it has almost invariant finite sets, F\o lner sets). Abelian and nilpotent groups are amenable; groups containing a free subgroup
are not.

\begin{think}
The number $\frac{\sqrt{2k-1}}k$ is $\frac{2\sqrt{d-1}}d$ with $d=2k$. The adjacency operator of the $d$-regular tree thus has norm $2\sqrt{d-1}$, although every vertex has degree $d$. Where does the walk ``lose'' the
missing mass, and why is a finite $d$-regular graph unable to do the same (Volume~11, Theorem~A.11)?
\end{think}

\section{Exercises}

\begin{exercise}[$\star$]\label{ex:dirichletprinciple}
Prove the Dirichlet principle: on a finite network, $\frac1{R(a\leftrightarrow Z)}=\min\{\mathcal E(f):f(a)=1,\ f=0\text{ on }Z\}$, attained at the voltage. Deduce Rayleigh's law again.
\end{exercise}

\begin{exercise}[$\star$]\label{ex:kac}
Deduce from the proof of Theorem~\ref{thm:commute} that $\mathbb E_a\tau_a^+=\frac{c(G)}{\pi(a)}$ (Kac's formula), and check it for the walk on a cycle and on a star.
\end{exercise}

\begin{exercise}[$\star$]\label{ex:hypercube}
On the hypercube $\{0,1\}^n$ with unit conductances, use symmetry and shorting to show that $R(0\leftrightarrow\mathbf 1)=\sum_{j=0}^{n-1}\frac1{(n-j)\binom nj}$, that $nR(0\leftrightarrow\mathbf 1)\to2$, and deduce that the
commute time between opposite corners is asymptotic to $2^{n+1}$.
\end{exercise}

\begin{exercise}[$\star\star$]\label{ex:spherical}
A \emph{spherically symmetric tree} has a root, and every vertex at distance $n$ from the root has $b_{n+1}\ge1$ children. Show that simple random walk on it is transient if and only if $\sum_n\frac1{b_1b_2\cdots b_n}<\infty$. Deduce that on a recurrent such tree the generation sizes $|T_n|=b_1\cdots b_n$
cannot all exceed $n^{1+\delta}$, and yet construct a recurrent one with $\limsup_n|T_n|/n^k=\infty$ for every $k$.
\end{exercise}

\begin{exercise}[$\star\star$]\label{ex:slab}
Show that simple random walk on the slab $\Z^2\times\{0,1,\dots,k\}$ is recurrent, and that on the triangular lattice ($\Z^2$ with the edges from $x$ to $x+(1,1)$ added) it is recurrent.
\end{exercise}

\begin{exercise}[$\star\star$]\label{ex:znreturn}
Show that $p_{2n}=\binom{2n}n^24^{-2n}$ for simple random walk on $\Z^2$ (rotate by $45^\circ$), and that $p_{2n}\le Cn^{-d/2}$ on $\Z^d$, so that $\rho=1$ on $\Z^d$. Compare with $F_2$, where $\rho=\frac{\sqrt3}2$: after $100$ steps,
how many orders of magnitude separate the return probabilities?
\end{exercise}

\begin{exercise}[$\star\star$]\label{ex:regulartree}
Let $\Gamma=\Z/2*\Z/2*\Z/2$, generated by three involutions $a,b,c$ with no other relation. Show that its Cayley graph is the $3$-regular tree and that $\rho=\frac{2\sqrt2}3$. More generally show that simple random walk on the
$d$-regular tree has $\rho=\frac{2\sqrt{d-1}}d$.
\end{exercise}

\begin{exercise}[$\star\star\star$]\label{ex:speed}
On the free group $F_k$, $k\ge2$, show that $|X_n|/n\to\frac{k-1}k$ almost surely, so the walk escapes at linear speed. Show that $X_n$ converges to an infinite reduced word $X_\infty$, and that the law of $X_\infty$ is the
uniform measure on the boundary of the tree (each first letter has probability $\frac1{2k}$, each next letter is uniform among the $2k-1$ allowed ones).
\end{exercise}

\section{Hints}

\begin{hint}{ex:dirichletprinciple}
Write $f=v+g$ with $g(a)=0$, $g=0$ on $Z$, and use \eqref{eq:green}.
\end{hint}

\begin{hint}{ex:kac}
In the proof, $\Delta\varphi_a(a)=\pi(a)-c(G)$; express $\Delta\varphi_a(a)$ with $\mathbb E_a\tau_a^+=1+\sum_yP(a,y)\varphi_a(y)$.
\end{hint}

\begin{hint}{ex:hypercube}
The voltage depends only on the number of ones. Short each level. The $j$-th level has $\binom nj$ vertices, each with $n-j$ edges to the next level. In the sum, the terms $j=0$ and $j=n-1$ dominate.
\end{hint}

\begin{hint}{ex:spherical}
Short the generations, as in Proposition~\ref{prop:treetransient}.
\end{hint}

\begin{hint}{ex:slab}
For the slab, cutsets made of the edges leaving the boxes $[-k',k']^2\times\{0,\dots,k\}$. For the triangular lattice, the same boxes as in $\Z^2$.
\end{hint}

\begin{hint}{ex:znreturn}
In $\Z^2$, $X_n=(U_n+V_n,U_n-V_n)/2$ with independent walks $U,V$ on $\Z$ of steps $\pm1$. For $\Z^d$, a local limit bound, or the Fourier formula $p_{2n}=(2\pi)^{-d}\int\bigl(\frac1d\sum\cos\theta_j\bigr)^{2n}\mathrm d\theta$.
\end{hint}

\begin{hint}{ex:regulartree}
Reduced words alternate letters. Repeat the proof of Theorem~\ref{thm:kesten} with $p=\frac{d-1}d$, $q=\frac1d$ and $f=(d-1)^{-|g|/2}$.
\end{hint}

\begin{hint}{ex:speed}
$|X_n|$ is a walk on $\{0,1,\dots\}$ with drift $p-q=\frac{k-1}k$ away from $0$, reflected at $0$; it visits $0$ finitely often. Once $|X_n|$ is large, the first letters are frozen.
\end{hint}

\section{Solutions}

\begin{solution}{ex:dirichletprinciple}
\idea Flows and potentials are dual: the current minimises energy among flows, and the voltage minimises energy among potentials. The two minima are $R$ and $1/R$.
\writeup Let $v$ be the voltage and $f=v+g$ with $g(a)=0$, $g|_Z=0$. By \eqref{eq:green}, $\mathcal E(f)=\mathcal E(v)+\mathcal E(g)+2\sum_xg(x)\Delta v(x)$, and the last sum vanishes because $\Delta v=0$ where $g\ne0$. So
$\mathcal E(f)\ge\mathcal E(v)$, with equality iff $\mathcal E(g)=0$, that is $g$ constant, hence $g=0$ since $g(a)=0$ and $G$ is connected. Finally $\mathcal E(v)=\sum_xv(x)\Delta v(x)=v(a)I=I=\frac1R$. When conductances increase, each $\mathcal E(f)$
increases, so $\frac1R$ increases.
\end{solution}

\begin{solution}{ex:kac}
\idea The walk from $a$ spends, between two visits to $a$, on average $c(G)/\pi(a)$ steps: the inverse of the stationary probability of $a$.
\writeup Since $\varphi_a(a)=0$ and $\mathbb E_a\tau_a^+=1+\sum_yP(a,y)\varphi_a(y)$, we have $\Delta\varphi_a(a)=-\pi(a)(\mathbb E_a\tau_a^+-1)$, and the proof showed $\Delta\varphi_a(a)=\pi(a)-c(G)$. Hence $\mathbb E_a\tau_a^+=\frac{c(G)}{\pi(a)}$. On a cycle of length $n$,
$c(G)=2n$ and $\pi(a)=2$: return time $n$. On a star with $m$ leaves, from the centre: $c(G)=2m$, $\pi=m$, return time $2$ (out and back); from a leaf: $\pi=1$, return time $2m$.
\end{solution}

\begin{solution}{ex:hypercube}
\idea Symmetry makes each level an equipotential, so shorting levels changes nothing and the cube becomes a path of bundles of parallel resistors.
\writeup The automorphisms of the cube fixing $0$ and $\mathbf 1$ act transitively on each level (permutations of coordinates), and they preserve the voltage, so it is constant on levels and no current flows between vertices
of the same potential when they are glued. The bundle between levels $j$ and $j+1$ has $(n-j)\binom nj$ parallel unit resistors. So $R=\sum_{j=0}^{n-1}\frac1{(n-j)\binom nj}$. The terms $j=0$ and $j=n-1$ are both $\frac1n$, the terms
$j=1$ and $j=n-2$ are $O(n^{-2})$, and the others are $O(n^{-3})$ each, at most $n$ of them: so $nR\to2$. By Theorem~\ref{thm:commute}, with $c(G)=2\cdot n2^{n-1}=n2^n$, the commute time is $n2^nR\sim2^{n+1}$.
\end{solution}

\begin{solution}{ex:spherical}
\idea Shorting generations is free by symmetry, and the tree becomes a path whose $n$-th edge is a bundle of $b_1\cdots b_n$ parallel wires.
\writeup The voltage from the root to the generation $N$ (shorted) depends only on the generation, by symmetry, so shorting each generation does not change the resistance. Between generations $n-1$ and $n$ there are
$b_1\cdots b_n$ edges in parallel, so $R(\text{root}\leftrightarrow\text{generation }N)=\sum_{n\le N}\frac1{b_1\cdots b_n}$, and $R(\text{root}\leftrightarrow\infty)$ is the full sum. Theorem~\ref{thm:recurrentR} concludes. If $|T_n|\ge n^{1+\delta}$ for large $n$, the series converges and the walk is transient. For the construction, let
$n_1=1$, $b_{n_j}=n_j^{\,j}$, $n_{j+1}=n_j+|T_{n_j}|$, and $b_n=1$ for the other $n$. Then $|T_{n_j}|\ge n_j^{\,j}$, so $|T_n|/n^k$ is unbounded for every $k$; but $|T_n|=|T_{n_j}|$ for $n_j\le n<n_{j+1}$, so each stretch contributes
$\frac{n_{j+1}-n_j}{|T_{n_j}|}=1$ to $\sum_n\frac1{|T_n|}$, which diverges: the walk is recurrent.
\end{solution}

\begin{solution}{ex:slab}
\idea Thickening the plane by a bounded amount multiplies the size of each cutset by a constant, which does not change the divergence of $\sum\frac1k$.
\writeup The edges joining $[-m,m]^2\times\{0,\dots,k\}$ to its complement form disjoint cutsets with $(k+1)(8m+4)$ edges, and $\sum_m\frac1{(k+1)(8m+4)}=\infty$: recurrent by Theorem~\ref{thm:nashwilliams}. In the triangular lattice the
edges leaving $[-m,m]^2$ are at most $3(8m+4)$, and the same argument applies.
\end{solution}

\begin{solution}{ex:znreturn}
\idea In two dimensions the walk is a pair of independent one-dimensional walks in rotated coordinates; in $d$ dimensions the Fourier integral is concentrated near $\theta=0$, where the characteristic function is $1-\frac{|\theta|^2}{2d}$.
\writeup With $U=X^1+X^2$ and $V=X^1-X^2$, each step changes $U$ and $V$ by $\pm1$, independently and uniformly, so $X_{2n}=0$ iff $U_{2n}=V_{2n}=0$, of probability $(\binom{2n}n4^{-n})^2$. For $\Z^d$, $p_{2n}=(2\pi)^{-d}\int_{[-\pi,\pi]^d}\phi(\theta)^{2n}\dd\theta$
with $\phi=\frac1d\sum\cos\theta_j$; near $0$, $|\phi|\le1-c|\theta|^2$, and near the points where $|\phi|=1$ (namely $\theta=0$ and $\theta=(\pi,\dots,\pi)$) the same bound holds after translation; elsewhere $|\phi|\le1-c'$. So
$p_{2n}\le C\int e^{-2nc|\theta|^2}\dd\theta+e^{-2nc'}(2\pi)^d\le C'n^{-d/2}$, and $p_{2n}^{1/(2n)}\to1$ because of the lower bound $p_{2n}=\sum_xP(X_n=x)P(X_n=-x)=\sum_xP(X_n=x)^2\ge(2n+1)^{-d}$ (symmetry and Cauchy--Schwarz, since $X_n$ takes at most $(2n+1)^d$ values). On $F_2$, $p_{100}\approx(\frac{\sqrt3}2)^{100}\approx5.7\times10^{-7}$ up to a polynomial factor, while on $\Z^2$, $p_{100}\approx\frac1{50\pi}\approx6\times10^{-3}$: four orders of magnitude, and the gap grows exponentially.
\end{solution}

\begin{solution}{ex:regulartree}
\idea Only the local structure matters: from every vertex but the start, one neighbour is closer and $d-1$ are farther.
\writeup Reduced words in $a,b,c$ are those with no two equal consecutive letters, and every element is a unique reduced word; from a reduced word one can append any of the $2$ letters different from the last, or cancel the
last: the Cayley graph is the $3$-regular tree. On the $d$-regular tree, Schur's test with $f=(d-1)^{-|g|/2}$ gives $\sum_yP(g,y)f(y)=\frac1d(\sqrt{d-1}+\frac{d-1}{\sqrt{d-1}})f(g)=\frac{2\sqrt{d-1}}df(g)$, and the distance chain with
$p=\frac{d-1}d$, $q=\frac1d$ gives the matching lower bound $\sqrt{4pq}=\frac{2\sqrt{d-1}}d$. For $d=3$, $\rho=\frac{2\sqrt2}3$.
\end{solution}

\begin{solution}{ex:speed}
\idea The distance from the identity is a biased walk; once far away, the walk never comes back to undo its first letters.
\writeup $|X_n|$ is the chain of the proof of Theorem~\ref{thm:kesten}. Off $0$ its increments are i.i.d.\ $\pm1$ with mean $p-q=\frac{2k-2}{2k}=\frac{k-1}k$; it is transient (Proposition~\ref{prop:treetransient} or $\rho<1$), so after its last visit to $0$
$|X_n|$ differs from a sum of $n$ i.i.d.\ increments of mean $\frac{k-1}k$ by at most twice the (finite) number of visits to $0$, and the strong law of large numbers gives $|X_n|/n\to\frac{k-1}k$. Since the biased walk on $\{1,2,\dots\}$ from $j$ returns to $j-1$ with probability $q/p<1$, each
level is visited finitely often, so for each $m$ the length exceeds $m$ from some time on; at a time after which $|X_n|>m$ forever, the first $m$ letters of the reduced word can no longer change (changing the $m$-th
letter requires going back to length $m-1$). So $X_n$ converges letter by letter to an infinite word $X_\infty$. By symmetry of the generators under the automorphisms of $F_k$ permuting $\{a_i^{\pm1}\}$, the first letter is
uniform; given the first $j$ letters, the automorphisms fixing them (and the subtree they define) permute the $2k-1$ allowed continuations, so the next letter is uniform among them.
\end{solution}

\section*{Sources and further reading}

\begin{itemize}
\item \emph{History.} G.~Kirchhoff, ``\"Uber die Aufl\"osung der Gleichungen, auf welche man bei der Untersuchung der linearen Vertheilung galvanischer Str\"ome gef\"uhrt wird'', \emph{Annalen der Physik und Chemie} 72 (1847),
497--508. C.~St.~J.~A.~Nash-Williams, ``Random walk and electric currents in networks'', \emph{Proceedings of the Cambridge Philosophical Society} 55 (1959), 181--194. H.~Kesten, ``Symmetric random walks on groups'',
\emph{Transactions of the American Mathematical Society} 92 (1959), 336--354, and ``Full Banach mean values on countable groups'', \emph{Mathematica Scandinavica} 7 (1959), 146--156. T.~Lyons, ``A simple criterion for
transience of a reversible Markov chain'', \emph{Annals of Probability} 11 (1983), 393--402. A.~K.~Chandra, P.~Raghavan, W.~L.~Ruzzo, R.~Smolensky and P.~Tiwari, ``The electrical resistance of a graph captures its commute
and cover times'', \emph{Proceedings of the 21st ACM Symposium on Theory of Computing} (1989), 574--586.
\item P.~G.~Doyle and J.~L.~Snell, \emph{Random Walks and Electric Networks}, Mathematical Association of America, 1984 (free at arXiv:math/0001057). The dictionary, with Rayleigh's and Thomson's principles and P\'olya's theorem.
\item R.~Lyons and Y.~Peres, \emph{Probability on Trees and Networks}, Cambridge University Press, 2016 (free online). Chapters~2, 3 and~6: flows, transience criteria, amenability and Kesten's theorem.
\item W.~Woess, \emph{Random Walks on Infinite Graphs and Groups}, Cambridge University Press, 2000. Spectral radius, amenability and the geometry of groups.
\end{itemize}
